\newcommand{\DoNotLoadEpstopdf}{}
\documentclass[11pt]{article}
\pdfinfoomitdate=1
\pdftrailerid{}
\usepackage[T1]{fontenc}
\usepackage{mathpazo}
\renewcommand{\ttdefault}{pcr}
\usepackage[margin=1in]{geometry}
\usepackage{amsmath,amssymb,amsthm,mathtools}
\usepackage{microtype}
\usepackage{enumitem,booktabs,array}
\usepackage{xcolor,tikz,tikz-cd}
\usetikzlibrary{arrows.meta,positioning,calc}
\usepackage[hidelinks]{hyperref}
\usepackage{xurl}

\newtheorem{theorem}{Theorem}[section]
\newtheorem{lemma}[theorem]{Lemma}
\newtheorem{proposition}[theorem]{Proposition}
\newtheorem{corollary}[theorem]{Corollary}
\theoremstyle{definition}
\newtheorem{definition}[theorem]{Definition}
\theoremstyle{remark}
\newtheorem{remark}[theorem]{Remark}
\newtheorem{notation}[theorem]{Notation}

\newcommand{\C}{\mathbb C}
\newcommand{\Q}{\mathbb Q}
\newcommand{\R}{\mathbb R}
\newcommand{\Z}{\mathbb Z}
\newcommand{\N}{\mathbb N}
\newcommand{\Gm}{\mathbb G_m}
\newcommand{\OO}{\mathcal O}
\newcommand{\PP}{\mathbb P}
\newcommand{\kk}{\mathbf k}
\DeclareMathOperator{\wt}{wt}
\DeclareMathOperator{\Spec}{Spec}
\DeclareMathOperator{\Supp}{Supp}
\DeclareMathOperator{\Div}{Div}
\DeclareMathOperator{\codim}{codim}
\DeclareMathOperator{\Hom}{Hom}
\DeclareMathOperator{\gr}{gr}
\DeclareMathOperator{\DR}{DR}
\DeclareMathOperator{\id}{id}
\DeclareMathOperator{\ord}{ord}
\DeclareMathOperator{\Exc}{Exc}
\DeclareMathOperator{\rk}{rk}
\DeclareMathOperator{\im}{im}
\DeclareMathOperator{\res}{res}
\DeclareMathOperator{\Sym}{Sym}
\DeclareMathOperator{\Frac}{Frac}
\DeclareMathOperator{\vol}{vol}
\DeclareMathOperator{\Pic}{Pic}
\DeclareMathOperator{\Cl}{Cl}
\DeclareMathOperator{\Bir}{Bir}
\DeclareMathOperator{\Aut}{Aut}
\DeclareMathOperator{\tr}{tr}
\DeclareMathOperator{\lct}{lct}
\DeclareMathOperator{\mult}{mult}
\DeclareMathOperator{\Gal}{Gal}
\DeclareMathOperator{\NS}{NS}

\setlist{itemsep=3pt,topsep=5pt}
\numberwithin{equation}{section}
\setcounter{tocdepth}{1}
\allowdisplaybreaks[1]
\raggedbottom
\hypersetup{bookmarksdepth=2,pdftitle={Relative denominators and effective systems for log Calabi-Yau fibrations},pdfauthor={OpenAI}}
\title{Relative denominators and effective systems for log Calabi-Yau fibrations}
\author{OpenAI}
\date{September 27, 2026}

\begin{document}
\maketitle
\begin{abstract}
We prove uniform denominator and effective-system bounds for log Calabi-Yau
fibrations from projective log canonical complex pairs of dimension at most
four onto positive-dimensional bases, with boundary coefficients in a fixed
finite rational set. The bounds give a uniform trivializing degree and a
uniform b-Cartier multiple of the moduli b-divisor. When the base divisor is
big, a uniform complete rounded adjoint system has section ratios generating
the full base function field.
\end{abstract}
\section{Introduction}\label{sec:introduction}

Let $(X,B)$ be a projective log canonical pair over $\C$, and let
$f:X\to Z$ be a contraction: a surjective projective morphism of normal
projective varieties with $f_*\OO_X=\OO_Z$.  We consider the case
\[
 K_X+B\sim_\Q f^*D,
\]
where $D$ is a rational Cartier divisor on $Z$.  The adjoint of the
generic fibre is then rationally linearly trivial.  The canonical bundle
formula separates the contribution of singular fibres from a moduli
part on the base.  Its qualitative form places that moduli part on a
birational model as a nef rational Cartier divisor.  For effective
applications, one needs a Cartier multiple controlled by the dimension
and boundary coefficients, rather than by the individual fibration.

This is a question about a chosen rational presentation.  Replacing a
base divisor by an arbitrary rationally linearly equivalent divisor
changes the moduli part by a rational principal divisor and can introduce
arbitrarily large denominators.  We prove that, when $\dim X\leq4$ and the
coefficients of $B$ belong to a fixed finite rational set, one can choose
an exact pullback presentation whose moduli part has a uniform Cartier denominator.

We recall the base data in the statement.  For the given $f$, the
\emph{discriminant} $B_Z$ has coefficient $1-t_P$ at a prime divisor
$P\subset Z$, where $t_P$ is the log canonical threshold of $f^*P$ over
the generic point of $P$.  This is computed on a neighbourhood of that
point where $P$ is Cartier.  The moduli data form a rational b-divisor
$\mathbf M$.  On a birational model $q:W\to Z$ where it descends, it is
represented by a rational Cartier divisor $M_W$ and its pullbacks to
higher models; its trace on $Z$ is $M_Z=q_*M_W$.  The b-divisor is
\emph{b-nef} if $M_W$ is nef, and $p\mathbf M$ is \emph{b-Cartier} if
$pM_W$ is Cartier.  A generalized pair $(Z,B_Z+M_Z)$ has
$K_Z+B_Z+M_Z$ rational Cartier, and its crepant boundaries are defined by
\[
 K_W+B_W+M_W=q^*(K_Z+B_Z+M_Z)
\]
and by the same rule on higher models.  It is generalized lc if all
coefficients of these boundaries are at most one, and generalized klt
if they are strictly less than one.  The crepant boundaries on higher
models may have negative coefficients.

\begin{theorem}[Uniform relative denominators]\label{thm:relative-denominators}
Fix a finite set $\Phi\subset[0,1]\cap\Q$ and an integer $1\leq d\leq4$.
There are positive integers $p_0\mid p$, with $p_0$ clearing $\Phi$, and a
rational DCC set $\mathcal B\subset[0,1]$, depending only on $d$ and
$\Phi$, with the following property. Suppose $(X,B)$ is a projective lc
pair of dimension $d$ with coefficients in $\Phi$, and
\[
 f:X\longrightarrow Z,\qquad \dim Z>0,\qquad K_X+B\sim_\Q f^*D,
\]
where $f$ is a contraction and $D$ is a rational Cartier divisor. There
exist $\psi\in\C(X)^*$ and $D_Z\sim_\Q D$ such that
\begin{equation}\label{eq:relative-exact-presentation}
 K_X+B+\frac{1}{p_0}\Div(\psi)=f^*D_Z,
 \qquad D_Z=K_Z+B_Z+M_Z.
\end{equation}
Here $B_Z$ is effective with coefficients in $\mathcal B$, the pair
$(Z,B_Z+M_Z)$ is generalized lc, and $\mathbf M$ is b-nef with
$p\mathbf M$ b-Cartier. If $(X,B)$ is klt, the generalized pair is
generalized klt.
\end{theorem}

The first equality in~\eqref{eq:relative-exact-presentation} is an
equality of actual rational divisors, for compatible choices of canonical
divisors and the displayed rational function.  It selects the
representative $D_Z\sim_\Q D$ to which the moduli bound applies.  The two
integers have distinct roles: $p_0$ gives the principal correction and,
on the generic fibre, a trivialization in that degree; $p$ clears the
moduli trace on a model determining $\mathbf M$.  The discriminant is
controlled instead by the fixed DCC set $\mathcal B$.

\subsection{Denominators, positivity, and earlier work}

Fujino and Mori bound the denominator of the semistable divisor
accompanying $bK_X$ in their canonical bundle formula, where $b$ is the
least nonvanishing pluricanonical degree of the generic fibre.  Their
bound uses the middle Betti number of a smooth model of a cyclic cover of
the fibre of Kodaira dimension zero.  The proof reduces to a curve and
controls a finite character through this cohomology
\cite[Theorem~3.1 and Sections~3.3, 3.6--3.8]{FujinoMori2000}.  In their klt log Iitaka
setting, Todorov and Xu obtain a bound for the moduli denominator in terms
of the boundary coefficients in relative dimension two, with arbitrary
total dimension \cite[Theorem~1.3]{TodorovXu2009}.  Their proof first
controls surface indices and the second Betti number of a resolution of
the cyclic cover, then applies the Fujino--Mori character method.

When the generic fibre is a rational curve, Floris bounds a common integer
clearing all moduli coefficients in terms of the fibre Cartier index.  Her
examples rule out every polynomial bound for such an integer as the index
varies \cite[Theorem~1.6]{Floris2013}.  For the lc
total spaces of dimension at most four considered here, the curve
reduction produces reduced special fibres of dimension at most three.
The index of the whole reduced fibre supplies the source of uniformity
in our proof.

Qualitative moduli theory supplies a different part of the argument.
Ambro establishes rational b-Cartier descent and positivity in the
klt-trivial setting \cite{AmbroBoundary2004,AmbroModuli2005}, and Fujino
and Gongyo extend b-nefness to lc-trivial fibrations
\cite[Theorem~3.6]{FujinoGongyoModuli}.  These results put the moduli data
on a suitable base model and make the determining divisor nef; their
Cartier multiple may depend on the fibration.  We use this descent and
nefness after fixing the exact rational trivialization and verifying the
required rank-one condition, and obtain a common multiple by a separate
special-fibre calculation.  The preprint
of Bakker, Filipazzi, Mauri and Tsimerman proves qualitative
b-semiampleness for lc-trivial fibrations with generically effective
boundary \cite[Theorem~1.5]{BakkerFilipazziMauriTsimerman2025}.  Their
Section~7.2.4 discusses the general effective question and known special
cases, where one seeks a uniform multiple that is base point free on a
determining model.  A bound for the Cartier denominator makes a multiple
Cartier on such a model; base point freeness is an additional
condition, and is not used here.

The argument on a reducible fibre also has an established ancestry.
Fujino's admissible and preadmissible sections organize pluricanonical
descent across the conductor
\cite[Definition~4.1 and Lemma~4.2]{Fujino2000}.  Gongyo uses this
framework in the abundance theorem for numerically trivial slc adjoints
\cite[Theorem~1.5 and Section~5]{Gongyo2013}.  Jiang and Liu supply the
uniform index theorem for projective slc log Calabi--Yau pairs through
dimension three that we use below
\cite[Corollary~1.6, arXiv version~1]{JiangLiu2021}.  The relevant output
is one trivialization on the slc pair itself, with its conductor gluing.
It is this global form that permits the cyclic-character comparison in
the present proof.

\subsection{The route to the relative bound}

Section~\ref{sec:generic} applies the lower-dimensional index theorem to
the geometric generic fibre, whose dimension is at most three, and
descends its principalization by Hilbert~90 in the same degree $p_0$.
This fixes the exact presentation~\eqref{eq:relative-exact-presentation}.
Threshold ACC supplies the DCC set for $B_Z$.  After applying the
qualitative canonical bundle formula and passing to a smooth determining
model $W\to Z$, it remains to control one moduli coefficient at a time.
For a prime $P\subset W$, let $\alpha$ be the coefficient of the pullback
of $D_Z$, and let $t_P$ be the threshold computed from the crepant sub-pair
on a resolved diagram over $W$.  The coefficient of $M_W$ differs from
$\alpha+t_P$ by an integer.  Section~\ref{sec:curve} takes a transverse
curve slice and runs a relative dlt MMP to obtain an exact local identity
on a dlt model.
Adjunction in Section~\ref{sec:residue} then gives an slc pair $(T,B_T)$
on the whole connected reduced special fibre, with $\dim T\leq3$ and
$K_T+B_T\sim_\Q0$.

The decisive uniform input is a form on this entire fibre.  Global ACC
places the different coefficients in a fixed finite set.  The Jiang--Liu
index theorem then gives a nowhere-vanishing log pluricanonical form $u$ on $T$ in one even
multiple $p$ of $p_0$.  The local residue test compares its pullback with
the residues of an ambient form on a normalized cyclic cover; that
ambient form carries the character recording the remaining denominator.
The comparison may use a larger, pair-dependent degree for ambient
adjunction, but it yields one scalar on each normalized component of the
connected reduced covering fibre $T_Y$.  In the uniform degree $p$, the
local conductor calculation shows that each tuple has matching
nonzero next residues at a generic crossing of two components.  Their
scalars therefore agree across intersections and hence on all of $T_Y$.
The ambient tuple is therefore one scalar multiple of the invariant
pullback of $u$, even
when the cyclic group permutes components.  Its character is trivial in
degree $p$, giving $p(\alpha+t_P)\in\Z$ and the required Cartier bound.

\subsection{Effective systems and torsion}

When $D$ is big, Proposition~\ref{prop:effective-base} gives an integer
$m=m(d,\Phi)$ such that, for every positive multiple $l$ of $m$, the
complete rounded divisorial system
\[
 \left|\left\lfloor l(K_X+B)\right\rfloor\right|
\]
is nonempty and its section ratios generate exactly
$\C(Z)\subset\C(X)$.  Here the sheaf is the rank-one reflexive divisorial
sheaf.  Equality of the embedded fields recovers the finite part of the
fibration as well as the dimension of its image.

For every positive $l$ divisible by $p_0$, Section~\ref{sec:effective}
uses the exact presentation to identify the full section spaces inside
$\C(X)$:
\[
 H^0\!\left(X,\OO_X(\lfloor l(K_X+B)\rfloor)\right)
 =\psi^{l/p_0}f^*H^0\!\left(Z,\OO_Z(\lfloor lD_Z\rfloor)\right),
\]
where $f^*$ means pullback of rational functions.  The valuation argument
works for reflexive sheaves without requiring $l(K_X+B)$ to be Cartier.
Birkar--Zhang's polarized effective birationality theorem applies to the
resulting big lc adjoint on $W$, with fixed DCC boundary coefficients and
$pM_W$ nef Cartier \cite[Theorem~1.3]{BirkarZhang2016}.  Choosing $m$
divisible by $p_0$ and their uniform degrees, the equality above transfers
the full base field to every required source system: the common factor
cancels in ratios.  Their separate Iitaka-fibration bound
\cite[Theorem~1.2]{BirkarZhang2016} retains dependence on a fibre
nonvanishing degree and a Betti number of a cyclic fibre cover.

Proposition~\ref{prop:rc-torsion} is a separate torsion result.  Fix
$1\leq e\leq4$, a rational DCC set $I\subset[0,1]$, and a positive
integer $p$.  For projective generalized klt pairs $(Z,B_Z+M_Z)$ of
dimension $e$, with boundary coefficients in $I$, b-nef rational data
$\mathbf M$ for which $p\mathbf M$ is b-Cartier, and a rationally
connected smooth projective resolution of $Z$, it gives one integer
$\ell=\ell(e,I,p)$ such that
\[
 K_Z+B_Z+M_Z\sim_\Q0
 \quad\Longrightarrow\quad
 \ell(K_Z+B_Z+M_Z)\text{ is an integral principal divisor}.
\]
Using suitable extractions, global ACC places the boundary coefficients
in a fixed finite set and gives a uniform positive lower bound for
generalized log discrepancies.  Birkar's boundedness theorem
\cite[Theorem~1.7]{BirkarBoundedCY} bounds the varieties up to isomorphism
in codimension one.  The proof combines their smooth-locus topology with
Kummer theory and finite generation of the divisor class group $\Cl(Z)$
from the rationally connected resolution to give a uniform exponent for
torsion in $\Cl(Z)$.
A separate uniform multiple clears the
coefficients of the actual adjoint before that torsion bound is applied.
Birkar's ordinary-pair index consequence, observed by Totaro
\cite[Corollary~1.8 and Section~9.5]{BirkarBoundedCY}, is a close predecessor.

Corollary~\ref{cor:relative-rc-torsion} combines this result with
Theorem~\ref{thm:relative-denominators}.  A projective klt fourfold pair with fixed finite
rational boundary coefficients and $K_X+B\sim_\Q0$ has a uniform integral
principal multiple of $K_X+B$ whenever it admits a contraction to a
positive-dimensional base with a rationally connected smooth projective
resolution.  The statement includes $B=0$.

\tableofcontents
\section{Divisors and the lower-dimensional index input}
\label{sec:preliminaries}

We fix the divisor conventions needed for the exact pullback identities,
the rounded systems, and the index theorem on the reduced fibre.
All varieties are over $\C$.  Unless stated otherwise, they are integral
and normal.  In the global statements, all varieties are projective,
including the bases of the contractions.  The reduced fibres may be
reducible and nonnormal, and their normalizations may have several
connected components.  In the local curve arguments, the total spaces
are projective over the indicated curves.  After a curve is shrunk to
an open neighbourhood, neither it nor its total space need be projective
over $\C$.  A contraction is a surjective projective morphism $f$ with
$f_*\OO=\OO$.

Boundaries are effective rational divisors; crepant sub-boundaries may
have negative coefficients.  We use the usual discrepancy conventions
for lc, klt, and dlt pairs.  A reduced pair $(T,B_T)$ is slc if $T$ is
pure-dimensional, $S_2$, and nodal in codimension one, no component of
$B_T$ is contained in the conductor, $K_T+B_T$ is rational Cartier, and
the normalization with the conductor included in the boundary is lc.  In
particular, the conductor has coefficient one on each normalized branch.
The fibres to which we apply this definition will be proved to have the
required pure dimension and depth properties.

Canonical divisors are chosen using rational top forms, and we retain
those choices in exact divisor identities.  For rational divisors on a
normal variety, $D_1\sim_\Q D_2$ means that some positive integer multiple of
$D_1-D_2$ is the divisor of a rational function.  An equality of
rational divisors is coefficientwise for the chosen representatives.
In particular, $rD$ being an \emph{integral principal
divisor} means
\[
 rD=\Div(v)
\]
for an actual rational function $v$; it includes integrality of every
coefficient of $rD$.

For an integral Weil divisor $G$ on a normal variety $X$, the reflexive
divisorial sheaf is
\[
 \OO_X(G)(U)=\{v\in\C(X):\Div(v)|_U+G|_U\geq0\}\cup\{0\}.
\]
For a rational divisor $D$, its rounded system is the complete system of
$\OO_X(\lfloor D\rfloor)$.  A nonzero rational function $v$ is a section
exactly when
\[
 \Div(v)+D\geq0,
\]
because its valuations are integers.  This criterion does not require
$D$ or $\lfloor D\rfloor$ to be Cartier.  The field generated by a
nonempty system is the subfield of $\C(X)$ generated by all ratios of
nonzero sections.  It distinguishes the full embedded field from a
proper subfield of the same transcendence degree.

A subset of $\R$ satisfies the descending chain condition (DCC) if it
contains no infinite strictly decreasing sequence.  The ascending chain
condition (ACC) is defined by reversing the inequalities.  Both enter
the proof: threshold ACC gives a DCC set for discriminant coefficients,
and global ACC will reduce certain numerically trivial pairs to finite
coefficient sets.

The index input needed for the relative theorem is in fibre dimension at
most three.  We use the following established theorem: for every finite
set $I\subset[0,1]\cap\Q$, there is a positive integer $n(I)$ such that
every projective slc pair $(T,B_T)$ of dimension at most three with
coefficients in $I$ and $K_T+B_T\sim_\Q0$ satisfies
\begin{equation}\label{eq:lower-slc-index}
 n(I)(K_T+B_T)\sim 0.
\end{equation}
This includes normal pairs and the zero boundary
\cite[Corollary~1.6, arXiv version~1]{JiangLiu2021}.  We may enlarge
$n(I)$ to clear $I$.  The conclusion means that the log pluricanonical
sheaf $\OO_T(n(I)(K_T+B_T))$ is invertible and trivial on $T$ itself.  A
choice of generator is a global form whose restrictions to the
normalized components satisfy the conductor gluing.  This global
trivialization, rather than separate trivializations on the components,
is the input required by the residue argument.

We also need this conclusion when the coefficients initially lie in a
fixed rational DCC set $I\subset[0,1]$.  On each normalized irreducible
component, include the conductor with coefficient one.  The resulting
pair is projective lc with effective boundary coefficients in the DCC
set $I\cup\{1\}$, and its adjoint is numerically trivial.  Global ACC
\cite[Theorem~1.5]{HMX2014} places all its nonzero coefficients in a
fixed finite rational subset, depending only on $I$ and the dimension
bound.  The coefficients of $B_T$ therefore also lie in a fixed finite
set.  Applying~\eqref{eq:lower-slc-index} to the slc pair $T$ supplies one
uniform global trivialization, including the gluing.  The possibly larger
degree used later to restrict an ambient pluricanonical sheaf to $T$ may
depend on the ambient pair; it is kept separate from this uniform index.

\section{An exact pullback and the coefficient to be controlled}
\label{sec:generic}

We begin the proof of Theorem~\ref{thm:relative-denominators}.  Put
$k=\dim Z$ and $K=\C(Z)$.  The contraction condition makes $K$
algebraically closed in $\C(X)$, so this function-field extension is
regular in characteristic zero.  The generic fibre is therefore
geometrically integral.  It is normal by localization from $X$, and
normality is geometric for a finite-type variety over the perfect field
$K$.  Choose a rational top form on $Z$.  Together with the form defining
$K_X$, the determinant sequence for function-field differentials determines
a rational relative top form on the generic fibre $X_\eta$.  We use its
divisor as $K_{X_\eta}$.  At codimension-one points of $X_\eta$, which are
smooth over $K$, this divisor is the coefficientwise localization of
$K_X$; vertical primes disappear.  If $d=k$, then $X_\eta=\Spec K$ and
the relative form is a nonzero $0$-form with zero divisor.
Restrict a log resolution and its crepant sub-boundary to the
geometric generic fibre.  Generic smoothness makes the horizontal strata
SNC there, so the original geometric generic pair is lc.  That pair is
projective of dimension $d-k\leq3$, its boundary coefficients belong to
$\Phi$, and its adjoint is rationally linearly trivial.

The lower-dimensional index theorem applies to this pair over
$\overline K$.  For completeness, $K$ is finitely generated over $\C$,
so $\overline K$ and $\C$ are algebraically closed fields of
characteristic zero with the same transcendence degree over $\Q$.
An abstract field isomorphism therefore identifies the pair with a
complex pair to which the theorem applies.

Choose a common integer $p_0$, also clearing $\Phi$, which principalizes
the adjoint of every such geometric generic fibre.  This
principalization descends to $K$ without enlarging the degree.  Indeed,
choose a finite Galois extension $L/K$ over which a principalizing
function is defined.  The quotient of any two of its Galois conjugates
has zero divisor on the proper geometrically integral normal fibre over
$L$, hence belongs to $L^*$.  These quotients form a multiplicative
Galois cocycle.  Hilbert~90 rescales the function to make it invariant,
and it then descends to $K$.  Consequently there is
$\psi\in\C(X)^*$ such that
\[
 p_0(K_X+B)+\Div(\psi)
\]
is vertical over $Z$.

Choose $a$ divisible by $p_0$ and $\phi\in\C(X)^*$ with
\[
 a(K_X+B)+\Div(\phi)=f^*(aD).
\]
The divisor of $\psi^{a/p_0}/\phi$ is vertical.  On the proper normal
generic fibre this function has neither zeros nor poles, so it belongs
to $K^*$.  Denoting it by $h$, set
\[
 D_Z=D+\frac1a\Div(h).
\]
Substitution gives the first equality
of~\eqref{eq:relative-exact-presentation} as an equality of actual
rational divisors.  In particular, $D_Z$ is rational Cartier and
$D_Z\sim_\Q D$.

We can now apply the lc-trivial-fibration theorem
\cite[Theorem~3.6]{FujinoGongyoModuli}.  Besides generic log canonicity
and the rational pullback relation, it requires a rank-one condition.
On a log resolution, the modified discrepancy divisor $\mathbf A^*$
omits the discrepancy $-1$ terms.  Because $0\leq B\leq1$,
rounding up leaves, over the generic point of the base, only an
effective exceptional divisor over $X$: strict transforms contribute
zero.  The associated divisorial sheaf pushes forward to $\OO_X$ by
normality.  Pushing on to $Z$ thus has rank one, as required.

For a prime divisor $P$ on any birational base model, let $t_P$ be the lc
threshold of its pullback over the generic point of $P$, computed with
the crepant sub-pair on a resolved diagram.  The discriminant
coefficient is $1-t_P$.  The crepant sub-pair is sub-lc, so $t_P\geq0$;
in the klt case it is sub-klt and $t_P>0$.  These inequalities on every base model give
the asserted generalized singularities.  On the original base $Z$,
the boundary upstairs is effective and a pullback component
dominating $P$ has positive integral multiplicity.  Hence $t_P\leq1$,
which gives $B_Z\geq0$.

For primes $P$ on $Z$, the thresholds satisfy ACC uniformly.  Shrink
about the generic point of $P$ so that $P$ is Cartier.  In a log
resolution computing the threshold, only divisors with positive order
along the pullback of $P$ matter.  Remove the images of those whose
images are proper subsets of $P$.  Then $t_P$ is an ordinary lc threshold of an effective
Cartier divisor, with positive integral coefficients, against an lc
pair with coefficients in $\Phi$.  The threshold ACC theorem
\cite[Theorem~1.1]{HMX2014} applies.  These thresholds are rational
and lie in $[0,1]$, so their complements form the required rational
DCC set $\mathcal B\subset[0,1]$.

The qualitative theorem makes the moduli b-divisor b-nef and gives its
descent to a nef rational Cartier divisor on a sufficiently high smooth
projective model $W\to Z$.  This holds for our chosen presentation:
changing a rational trivialization adds a rational principal b-divisor
from the base, which preserves nefness and rational Cartier descent.
Write
\[
 D_W=(W\to Z)^*D_Z,
 \qquad M_W=D_W-K_W-B_W.
\]
For $P\subset W$, put $\alpha=\operatorname{coeff}_P D_W$.  Since
$\operatorname{coeff}_P B_W=1-t_P$ and the coefficient of $K_W$ is an
integer,
\[
 \operatorname{coeff}_P M_W
   =\alpha+t_P-1-\operatorname{coeff}_P K_W.
\]
It remains to prove, uniformly in $P$, that
\begin{equation}\label{eq:relative-coefficient-target}
 p(\alpha+t_P)\in\Z.
\end{equation}
Then $pM_W$ is an integral divisor on the smooth variety $W$, hence Cartier.

\section{Reduction to a dlt fibre over a curve}
\label{sec:curve}

Our objective is the integrality condition~\eqref{eq:relative-coefficient-target}.
Fix a prime $P\subset W$.  Take a smooth projective model $V\to X$ mapping
to $W$ by $g$, with SNC markings for the crepant boundary $B_V^c$, the
exceptional divisors over $X$, and the support of $g^*P$.
The morphism $g$ has connected fibres: its function field extension is
$\C(W)=\C(Z)\subset\C(V)=\C(X)$, which is regular. Its Stein
factor is therefore finite birational over the normal variety $W$ and
hence is $W$ itself.  Let $\theta_V$ be the rational top form compatible
with the chosen canonical divisor on $X$, and put $K_V=\Div(\theta_V)$.
Crepant pullback gives
\begin{equation}\label{eq:relative-resolution-presentation}
 K_V+B_V^c+\frac1{p_0}\Div(\psi)=g^*D_W.
\end{equation}
Put $t=t_P$.  At the generic point of $P$, the boundary
$B_V^c+t g^*P$ has coefficients at most one, with equality on at least one
component dominating $P$.  This is the SNC threshold calculation.

Choose sufficiently general very ample divisors $A_1,\ldots,A_{k-1}$ on
$W$ and put $C=A_1\cap\cdots\cap A_{k-1}$, taking $C=W$ when $k=1$.
Bertini, applied successively to the pullback linear systems and the
finitely many marked strata, makes $C$ and $V_C=g^{-1}C$ smooth, with
SNC restricted markings.  The curve $C$ is connected, and the restriction
of $g$ has connected fibres, so $V_C$ is connected and hence integral.
We also require $C\not\subset\Supp D_W$ and
$V_C\not\subset\Supp K_V\cup\Supp B_V^c\cup\Supp\Div(\psi)$.
The dimension of $V_C$ is $d-k+1\leq4$.  Choose an intersection point
$c\in C\cap P$ transverse and general on $P$.  We may avoid at $c$ all
other components of $\Supp D_W$, the bad strata loci on $P$, and the
images of all relevant nondominating strata which do not contain $P$.

Let $h:V_C\to C$, $F=h^*[c]$, and $B^c=B_V^c|_{V_C}$.  We specify the
canonical divisor in complete-intersection adjunction in order to retain
the coefficient of $P$.  Choose general representatives $A_j'\sim A_j$
avoiding $c$, and rational functions $\chi_j\in\C(W)^*$ with
$\Div(\chi_j)=A_j'-A_j$.  Define $\theta_C$ as the iterated Poincar\'e
residue, in the order of the cuts, of
\[
 \theta_V\prod_{j=1}^{k-1}g^*\chi_j
\]
along $g^{-1}A_1,\ldots,g^{-1}A_{k-1}$, taking each residue on the
preceding cut.  The generality of the divisors makes this a nonzero
rational top form on $V_C$.  For $k=1$ the product is empty and the residue
is the identity.  Each $\chi_j$ supplies a simple pole along the cutting
divisor and a zero along its chosen representative.  Cancelling the cutting
divisors before restriction, the residue formula gives an equality of
actual divisors:
\[
\begin{aligned}
 K_{V_C}:=\Div(\theta_C)
     &=\left(\Div\left(\theta_V\prod_jg^*\chi_j\right)
                   +\sum_jg^*A_j\right)|_{V_C}\\
     &=\left(K_V+\sum_jg^*A_j'\right)|_{V_C}.
\end{aligned}
\]
Set $\psi_C=\psi|_{V_C}$ and $D_C'=(D_W+\sum_jA_j')|_C$.  All these
restrictions are defined by generality.  Restricting
\eqref{eq:relative-resolution-presentation} now gives
\begin{equation}\label{eq:relative-curve-presentation}
 K_{V_C}+B^c+\frac1{p_0}\Div(\psi_C)=h^*D'_C,
 \qquad \operatorname{coeff}_c D'_C=\alpha.
\end{equation}
The coefficient assertion follows because $C$ meets $P$ transversely at
$c$, while the $A_j'$ and all other components of $\Supp D_W$ avoid $c$.

Along $F$, the coefficients of $B^c+tF$ are at most one, with equality on a
component of $F$.  Indeed, the marked divisors dominating $P$ restrict
transversely and reducedly, and near the selected point
$g^*P|_{V_C}=F$.  Over the generic point of $C$ there is still a projective
birational comparison with the normal fibre of the original model $X$,
and the restricted exceptional divisors are exceptional there.  To justify
this last assertion, the generic point of $C$ lies where $W\to Z$ is an
isomorphism.  Generic flatness and openness of geometric normality and
integrality give normal integral original fibres on an open base, and the
birational isomorphism locus is fibrewise dense there.  Shrinking further
makes the exceptional images have fibre codimension at least two.

On $V_C$ define an effective boundary $\Gamma$ by taking $F_{\mathrm{red}}$,
the horizontal strict transforms of the original boundary with their
original coefficients, and all horizontal restricted exceptional divisors
with coefficient one.  Then $(V_C,\Gamma)$ is log smooth, its horizontal
coefficients belong to the fixed set $\Phi\cup\{1\}$, and
\begin{equation}\label{eq:relative-error}
 E=\Gamma-B^c-tF
\end{equation}
is effective near $c$ and has coefficient zero on a component of $F$.
On the generic fibre $E$ is effective: its horizontal nonexceptional
terms cancel, while every horizontal exceptional term has coefficient
$1-\operatorname{coeff}B^c\geq0$.  Its support there is exceptional over
the original normal fibre by the preceding birational comparison.
Equation~\eqref{eq:relative-curve-presentation} gives the exact identity
\[
 K_{V_C}+\Gamma+\frac1{p_0}\Div(\psi_C)
     =h^*(D'_C+t[c])+E.
\]

We run a minimal-model program over $C$.  All negative coefficients of
$E$ are vertical and occur over finitely many points, so there is an
effective rational divisor $A$ on $C$ with $E^+:=E+h^*A\geq0$.  The
preceding identity becomes
\[
 K_{V_C}+\Gamma+\frac1{p_0}\Div(\psi_C)
     =h^*(D'_C+t[c]-A)+E^+.
\]
In particular, $K_{V_C}+\Gamma\sim_{\Q,C}E^+$; its restriction to a very
general fibre is rationally linearly equivalent to an effective divisor,
so it is pseudo-effective over $C$.

The existence and preservation results recalled in
\cite[Remark~2.11, arXiv version~2]{ChenTsakanikas2023} permit an MMP with
scaling of an ample divisor over the base for an lc generalized pair whose
underlying variety is $\Q$-factorial klt.  We apply this to the
$\Q$-factorial dlt pair $(V_C,\Gamma)$ over $C$; the same reference states
that its stages remain $\Q$-factorial and dlt.  Each birational step is
negative for the transformed adjoint.  The zero nef datum satisfies the
source's NQC hypothesis.  In dimension four,
\cite[Theorem~1.1, arXiv version~2]{ChenTsakanikas2023} terminates every
sequence of flips over $C$ for an NQC lc generalized pair whose adjoint is
pseudo-effective over $C$.  The three-dimensional statement
\cite[Theorem~1.2, arXiv version~2]{ChenTsakanikas2023} has no
pseudo-effectivity hypothesis.

These results supply the program and terminate a possible flip tail.
We check that its endpoint is nef and retains the exact identity.  Use
$\theta_C$ to choose canonical divisors on every birational stage
$h_i:V_i\to C$, and let $\Gamma_i$ and $E_i^+$ denote the pushforwards.
Every step is over $C$, so the preceding identity pushes forward to
\[
 K_{V_i}+\Gamma_i+\frac1{p_0}\Div(\psi_C)
     =h_i^*(D'_C+t[c]-A)+E_i^+,\qquad E_i^+\geq0.
\]
Thus the adjoint remains pseudo-effective over $C$.  A fibre-type negative
extremal contraction is impossible: a general positive-dimensional
contracted fibre has a curve $\ell$ not contained in $\Supp E_i^+$.
Since $V_i$ is $\Q$-factorial, a Cartier multiple of $E_i^+$ restricts
to an effective divisor on the normalization of $\ell$.  Hence
\[
 0\leq E_i^+\cdot\ell=(K_{V_i}+\Gamma_i)\cdot\ell<0,
\]
a contradiction.  Divisorial contractions strictly decrease the relative
Picard number, so an infinite program would eventually consist only of
flips.  The cited termination statements exclude such a tail in dimensions
three and four; in dimensions at most two no flips occur.  The endpoint is
a projective morphism $f_N:N\to C$ with $N$ $\Q$-factorial,
$(N,\Gamma_N)$ dlt, and $K_N+\Gamma_N$ nef over $C$.  Its unchanged regular
function-field extension and Stein factorization give connected fibres,
so $f_N$ is a contraction.

Let $E_N^+$ be the transform of $E^+$ on $N$.  The pushforward of
$h^*A$ is $f_N^*A$, so
$E_N:=E_N^+-f_N^*A$ is the pushforward of the original $E$.  Substituting
$E_N^+=E_N+f_N^*A$ into the transformed identity gives
\begin{equation}\label{eq:relative-output-error}
 K_N+\Gamma_N+\frac1{p_0}\Div(\psi_C)
   =f_N^*(D'_C+t[c])+E_N.
\end{equation}

We claim that $E_N=0$ near $c$.  Take a common resolution with maps
$r:\widetilde V\to V_C$ and $r':\widetilde V\to N$.  Discrepancy
comparison for the $K+\Gamma$ MMP gives
\[
 r^*(K_{V_C}+\Gamma)-r'^*(K_N+\Gamma_N)\geq0
\]
by \cite[Lemmas~3.38--3.39]{KollarMori1998}.  Pulling back the exact input
and output identities cancels the common base pullbacks and the common
principal divisor, giving
\begin{equation}\label{eq:relative-error-monotonicity}
 r'^*E_N\leq r^*E.
\end{equation}
Let $\eta$ be the generic point of $C$ and
$q:\widetilde V_\eta\to X_\eta$ the morphism to the normal original fibre.
The restriction of $E_N$ to $N_\eta$ is effective, being the pushforward
of the effective restriction of $E$; hence the restriction of $r'^*E_N$
is effective as well.  The restriction of $r^*E$ is effective and
$q$-exceptional, so \eqref{eq:relative-error-monotonicity} makes the
restriction of $r'^*E_N$ $q$-exceptional.  Finally,
\eqref{eq:relative-output-error} gives
$E_N\equiv_C K_N+\Gamma_N$, so that restriction is nef, in particular
$q$-nef.  The negativity lemma makes it zero.  Thus $E_N$ has no
horizontal component.  After shrinking about $c$, it is effective and
supported on the special fibre, and it remains nef over $C$.

Write
\[
 F_N=f_N^*[c]=\sum_i e_iS_i,\qquad E_N=\sum_i q_iS_i,
 \qquad e_i>0,\quad q_i\geq0,
\]
and set $\lambda=\max_i(q_i/e_i)$.  The divisor
$\Delta=\lambda F_N-E_N$ is effective, anti-nef over $C$, and misses at
least one component of the fibre.  Suppose the fibre has positive
dimension and $\Delta\ne0$.  Connectedness gives a component $S$ outside
$\Supp\Delta$ that meets it.  Since $N$ is $\Q$-factorial, choose a
positive integer $b$ so that $b\Delta$ is Cartier.  Its restriction to the
integral projective component $S$ is a nonzero effective Cartier divisor.
If $n=\dim S$ and $A_S$ is a very ample
Cartier divisor on $S$, then
\[
 \bigl((b\Delta)|_S\bigr)\cdot A_S^{n-1}>0.
\]
Here $A_S^{n-1}$ is represented by an effective curve cycle in the fibre
(by $S$ itself when $n=1$), so anti-nefness gives the opposite inequality.
This contradiction proves $\Delta=0$, hence $E_N=\lambda F_N$.
For a zero-dimensional fibre, $f_N$ is an isomorphism of normal curves
and the same proportionality is immediate.  Finally,
\[
 r'^*F_N=r^*F,
\]
and the original $E$ has coefficient zero on a component of $F$.
Taking the coefficient of its strict transform in
\eqref{eq:relative-error-monotonicity} forces $\lambda=0$, whether or not
that component survives on $N$.
This proves the claim.

Shrink $C$ about $c$ and set $T=(f_N^*[c])_{\mathrm{red}}$.  No divisor was
extracted in the program, so every surviving component of this fibre has
coefficient one in $\Gamma_N$.  Thus
$\Gamma_N=T+H$, where $H$ is horizontal and has coefficients in the fixed
set $\Phi\cup\{1\}$.  The reduced fibre $T$ is connected and projective.
Its components have dimension $d-k\leq3$, since they are the components
of the effective Cartier divisor $f_N^*[c]$ on the integral variety $N$.
Removing all boundary from the $\Q$-factorial
dlt pair shows that $N$ is klt.  Shrink again so that $D'_C$ is supported
at most at $c$.  For the integral Weil divisor
$L=p_0(K_N+T+H)$, the exact relation is now
\begin{equation}\label{eq:relative-local-cyclic}
 L+\Div(\psi_C)=p_0(\alpha+t)f_N^*[c].
\end{equation}

\section{The denominator detected by residues on the whole fibre}
\label{sec:residue}

In the application to \eqref{eq:relative-local-cyclic}, the coefficient
below is $\beta=\alpha+t$.  Write $p_0\beta=a/m$ in lowest terms, with
$m>0$.  An $m$th root of a base uniformizer produces a log
pluricanonical generator with character $\zeta^{-a}$ for
$\zeta\in\mu_m$.  We compare its
$p/p_0$th power with the pullback of a degree-$p$ form on the whole
reduced fibre.  Divisible adjunction will first make their ratios
constant on each normalized component.  Residues along the conductor
will then make those constants equal, forcing the character to be
trivial and hence $m\mid p/p_0$.

\begin{lemma}[The reduced fibre and the cyclic residue test]
\label{lem:slc-residue-test}
Fix a finite set $\Phi\subset[0,1]\cap\Q$ and a positive integer
$p_0$ clearing its denominators.  There is a positive integer $p$,
depending only on $\Phi$ and $p_0$, divisible by $p_0$, with the following
property.  Let
\[
 f\colon N\longrightarrow C
\]
be a projective contraction from a normal integral variety to a smooth
complex curve, let $c\in C$, and suppose $1\leq\dim N\leq4$.  Set
$F=f^*[c]$ and $T=F_{\mathrm{red}}$.  Assume that $N$ is
$\Q$-factorial, that $(N,T+H)$ is dlt, and that $H$ is an effective
horizontal $\Q$-divisor with coefficients in $\Phi$.  Suppose, on a
neighbourhood of $F$, that there are $\beta\in\Q$ and
$\psi\in\C(N)^*$ such that the following is an equality of actual
$\Q$-divisors, for a fixed canonical divisor:
\begin{equation}
 \label{eq:slc-local-relation}
 p_0(K_N+T+H)+\Div\psi
       =p_0\beta F.
\end{equation}
Then $p\beta\in\Z$.
\end{lemma}

\begin{proof}
We may shrink $C$ about $c$ throughout.  In particular, we take a
uniformizer $z\in\C(C)$ with $\Div_C z=[c]$ on this
neighbourhood.  The fibre $T$ is connected and projective, since $f$ is
a projective contraction.  Its irreducible components have dimension
$\dim N-1$: the divisor $F$ is an effective Cartier divisor on the
integral variety $N$, and its support is the whole fibre.  When
$\dim N=1$, the contraction $f$ is an isomorphism of normal curves.
Comparing the coefficient at $c$ in
\eqref{eq:slc-local-relation} gives $p_0\beta\in\Z$, so this case is
immediate.  Henceforth $1\leq\dim T\leq3$.

\smallskip
\noindent\emph{Adjunction on the entire reduced fibre.}
Write $T=\bigcup_i T_i$.  By dlt adjunction and the structure theorem for
dlt strata, each $T_i$ is normal and
\begin{equation}
 \label{eq:slc-component-adjunction}
 (K_N+T+H)|_{T_i}\sim_{\Q}K_{T_i}+\Theta_i,
\end{equation}
where $(T_i,\Theta_i)$ is lc and $\Theta_i$ is the effective different;
see \cite[Sections~4.1--4.2 and Theorems~4.16 and~4.19]{Kollar2013}.
We use the residue form of this adjunction
\cite[Definition~11.14, author version]{KollarFamilies2023}.  If $n$ is sufficiently
divisible, the $n$th tensor power of the ordinary residue at the generic
point of $T_i$ extends to an isomorphism
\[
 \OO_N\bigl(n(K_N+T+H)\bigr)|_{T_i}
   \simeq \OO_{T_i}\bigl(n(K_{T_i}+\Theta_i)\bigr).
\]
The different is defined so that the generic residue extends uniquely
to this canonical isomorphism once the ambient Cartier index is
cleared.  Here $T_i$ is normal, so its normalization does not change
the target.  The isomorphism therefore identifies actual rational
pluricanonical forms.  The same form of adjunction will be used for
normal dlt components upstairs.
Intersections of distinct components of $T$ are unions of lc strata.
At the generic point of an intersection of codimension two in $N$, the
pair is simple normal crossing, precisely two components of $T$ pass
through the point, and $H$ is absent.  The corresponding conductor
prime on either branch occurs in $\Theta_i$ with coefficient one.

The remaining coefficients of the differents belong to the set
\begin{equation}
 \label{eq:slc-different-coefficients}
 \mathcal D(\Phi)=
 \left\{
   \frac{r-1+b}{r}\ \middle|\
   \begin{gathered}
    r\in\N_{>0},\quad n_\phi\in\Z_{\geq0},\\
    b=\sum_{\phi\in\Phi\cup\{1\}}n_\phi\phi\leq1
   \end{gathered}
 \right\}\cap[0,1].
\end{equation}
This is the usual coefficient rule for the different; see
\cite[Lemma~4.1]{HMX2014}.  The positive elements of
$\Phi\cup\{1\}$ have a positive minimum, so the possible sums $b\leq1$
form a finite set.  For each such sum, the displayed coefficients are
nondecreasing in $r$ and converge to one.  Consequently $\mathcal D(\Phi)$ is a
fixed rational DCC set.  We include the conductor coefficient one in
this set.

To apply the index theorem to $T$, we must glue this component
adjunction on the reduced scheme itself.  We will prove that $T$ is
demi-normal and that divisible residues identify its log
pluricanonical sheaf with an ambient restriction.  The degree needed
for this adjunction may depend on the pair.  Conductor-compatible
descent is also the issue addressed by the admissible-section methods
of Fujino \cite[Definition~4.1 and Lemma~4.2]{Fujino2000} and Gongyo
\cite[Section~5]{Gongyo2013}; here we give the local comparison needed
for the later uniform degree.

Apply the depth theorem
\cite[Theorem~11.18, author version]{KollarFamilies2023} to the dlt pair
$(N,T+H)$ with the effective integral divisor
$D=T\leq\lfloor T+H\rfloor$ and $L=0$.  It gives that
$\OO_N$, $\OO_N(-T)$, and $\OO_T$ are Cohen--Macaulay.
Here $\OO_N(-T)$ is the ideal of the reduced union of the prime
divisors $T_i$: on the normal variety $N$, it is the intersection of
their height-one prime ideals, equivalently the regular functions
vanishing to order at least one at every $T_i$.  Thus the subscheme in
the depth theorem is the reduced fibre used here.  In particular, $T$
satisfies $S_2$.

The normality of the $T_i$ and the generic normal crossing description
show that $T$ has only smooth points and ordinary double crossings in
codimension one.  It is also seminormal.  By the square--cube
criterion, an element of the total quotient ring whose square and cube
are regular lies in the local ring at each codimension-one point,
where seminormality is explicit.  The $S_2$ extension property then
puts it in $\OO_T$.  Thus $T$ is demi-normal.  Its normalization is the
disjoint union of the $T_i$, and its conductor divisor is the union of
the double-crossing divisors described above.

Remove the coefficient-one conductor divisors from the $\Theta_i$ and
push the remaining boundary cycles to $T$; denote the result by $B_T$.
Off the conductor the normalization is an isomorphism in codimension
one, so this does not double count any prime.  The boundary $B_T$ is
effective, has coefficients in $\mathcal D(\Phi)$, and has no component
contained in the conductor.

Choose an open immersion $j:T^\circ\hookrightarrow T$ whose complement
has codimension at least two, such that $T^\circ$ has only smooth points
and ordinary double crossings.  We may also arrange that at each crossing
the ambient pair is simple normal crossing with just the two vertical
branches, and that the support of $B_T|_{T^\circ}$ is contained in the
smooth locus.  The dualizing
sheaf of $T^\circ$ is invertible, and the relevant multiples of
$B_T|_{T^\circ}$ are Cartier.  For such a multiple $n$, interpret
$\OO_T(n(K_T+B_T))$ initially as
\[
 j_*\bigl(\omega_{T^\circ}^{\otimes n}
       \otimes\OO_{T^\circ}(nB_T|_{T^\circ})\bigr).
\]
We claim that for every sufficiently divisible positive integer $n$,
the residue maps give a natural isomorphism
\begin{equation}
 \label{eq:slc-ambient-restriction}
 \OO_N\bigl(n(K_N+T+H)\bigr)|_T
       \simeq \OO_T\bigl(n(K_T+B_T)\bigr).
\end{equation}
Choose $n$ even and divisible enough that the ambient sheaf is
invertible and the component residue isomorphisms hold.  Away from the
conductor, these isomorphisms give
\eqref{eq:slc-ambient-restriction} on $T^\circ$.  At a node, take local
normal crossing coordinates $x,y$ for the two branches.  Hypersurface
adjunction for $xy=0$ identifies the restriction of the ambient log
canonical line with the dualizing line of the node.  The two iterated
residues of $dx/x\wedge dy/y$ differ by a sign; their $n$th tensor
powers give the dualizing-sheaf gluing and have equal next residues
when $n$ is even.  This proves the claimed residue identification on
$T^\circ$.

An invertible sheaf on the $S_2$ scheme $T$ equals the extension of
its restriction from $T^\circ$.  Applying this to the left side of
\eqref{eq:slc-ambient-restriction} and using the definition of the
right side proves the isomorphism on $T$.  In particular, the defined
extension is invertible in these degrees, and the isomorphism retains
the actual rational residue forms on every normalized branch.

It follows that $K_T+B_T$ is $\Q$-Cartier.  Its pullback to the
normalization is $K_{T_i}+\Theta_i$, with the conductor included, so
$(T,B_T)$ is slc.  Equation~\eqref{eq:slc-local-relation}, together with
$F=\Div(f^*z)$, makes $K_N+T+H$ rationally linearly trivial
near $T$.  Restricting a sufficiently divisible trivialization in
\eqref{eq:slc-ambient-restriction} gives
\begin{equation}
 \label{eq:slc-fibre-trivial}
 K_T+B_T\sim_{\Q}0.
\end{equation}

\smallskip
\noindent\emph{A uniform degree on the slc fibre.}
Apply global ACC to the projective normalized component pairs
$(T_i,\Theta_i)$, whose adjoints are numerically trivial by
\eqref{eq:slc-fibre-trivial}.  Their dimensions are at most three and
their coefficients lie in the fixed DCC set
$\mathcal D(\Phi)$; hence their nonzero coefficients lie in a fixed
finite rational set \cite[Theorem~1.5]{HMX2014}.  The same is true of
$B_T$.  The bounded index theorem for projective slc log Calabi--Yau
pairs in dimensions at most three now gives a common integer
annihilating $K_T+B_T$ \cite[Corollary~1.6, arXiv version~1]{JiangLiu2021}.
Enlarging this integer, choose once and for all an even multiple $p$ of
$p_0$, depending only on $\Phi$ and $p_0$, such that
\begin{equation}
 \label{eq:slc-uniform-index}
 pB_T\text{ is integral},\qquad p(K_T+B_T)\sim0.
\end{equation}
Choose a nowhere-vanishing generator $u$ of this trivial log
pluricanonical line.  On the normalization it is a tuple of rational
$p$-pluricanonical forms $u_i$ on the $T_i$, with the prescribed boundary
poles and conductor gluing.

The comparison uses two degrees for different purposes.  A sufficiently
divisible degree $n=bp$ makes $u^b$ a generating ambient restriction
and will control divisors of componentwise ratios.  The conductor
comparison will then take place in the uniform degree $p$, using only
the ordinary node description there.

\smallskip
\noindent\emph{The root cover.}
Write
\begin{equation}
 \label{eq:slc-denominator}
 p_0\beta=\frac{a}{m},\qquad
 a\in\Z,\quad m\in\N_{>0},\quad\gcd(a,m)=1.
\end{equation}
Thus the goal is $m\mid p/p_0$. The case $m=1$ is immediate; assume $m>1$.  Let $C'\to C$ be the normalization in
$\C(C)(w)$, where $w^m=z$, and let $\pi\colon Y\to N$ be the
normalization of the dominating component of $N\times_C C'$.  The
polynomial $w^m-z$ is Eisenstein at $c$.  Thus $C'$ has a unique point
$c'$ over $c$, with total ramification of degree $m$ there.  Since $f$
is a contraction of normal varieties, $\C(C)$ is algebraically closed
in $\C(N)$; in characteristic zero this is a regular function-field
extension.  It is consequently linearly disjoint from
$\C(C')/\C(C)$.  The map $\pi$ has degree $m$ and cyclic Galois group
$\mu_m$, acting by $w\mapsto\zeta w$.
Write $f_Y:Y\to C'$ for the induced morphism and put
$T_Y=(\pi^{-1}T)_{\mathrm{red}}$. Figure~\ref{fig:root-cover}
records these maps and their reduced special fibres; the total space
$Y$ is the normalization of the base change.

\begin{figure}[ht]
\centering
\begin{tikzcd}[row sep=large,column sep=huge]
 T_Y \arrow[r] \arrow[d,hook] & T \arrow[d,hook] \\
 Y \arrow[r,"\pi"] \arrow[d,"f_Y"'] & N \arrow[d,"f"] \\
 C' \arrow[r,"w^m=z"'] & C
\end{tikzcd}
\caption{The normalized cyclic base change and its whole reduced special
fibres over $c'$ and $c$. The covering group may permute components of
$T_Y$, so the residue comparison must use their conductor gluing.}
\label{fig:root-cover}
\end{figure}

The morphism $f_Y$ is projective.  Its function
field is regular over $\C(C')$, because regularity is preserved by this
base change.  Stein factorization therefore gives
$(f_Y)_*\OO_Y=\OO_{C'}$. In particular, $T_Y$ is connected and
projective.
The divisor $p_0(K_N+T+H)$ in
\eqref{eq:slc-local-relation} is integral.  If
$F=\sum_i e_iT_i$, comparison of coefficients in that equation and
\eqref{eq:slc-denominator} gives $m\mid e_i$ for every $i$.
At the generic point of $T_i$, write $f^*z=\tau^{e_i}v$ with $v$ a unit.
After dividing $w$ by $\tau^{e_i/m}$, the normalized extension is
obtained by extracting an $m$th root of $v$ and is \'{e}tale.  Away from
$T$ the function $f^*z$ is a unit, and the same conclusion holds.
Hence $\pi$ is \'{e}tale in codimension one.

Put $H_Y=\pi^*H$, and choose canonical divisors compatibly.  Since
$\pi$ is quasi-\'{e}tale and the coefficient-one divisor pulls back
without divisorial ramification, we have
\begin{equation}
 \label{eq:slc-cover-crepant}
 K_Y+T_Y+H_Y=\pi^*(K_N+T+H).
\end{equation}
The pair on the left is dlt.  Indeed, finite crepant discrepancy
comparison multiplies a downstairs log discrepancy by the
ramification index of a divisorial prolongation
\cite[Proposition~5.20]{KollarMori1998}.  Thus the upstairs pair is lc,
and every upstairs lc centre maps onto a downstairs lc centre.
The pair downstairs is simple normal crossing near the generic point
of each such centre.  A finite quasi-\'{e}tale map is \'{e}tale over the
smooth locus of its target, by purity of the branch locus
\cite[Tag~0BMB]{StacksProject}.  Consequently the upstairs pair is
simple normal crossing at the generic points of all its lc centres.
A log resolution chosen to be an isomorphism over this log smooth open
has no exceptional divisor of log discrepancy zero, which is the dlt
condition.  This argument does not require $Y$ to be $\Q$-factorial.
In particular, each irreducible component $T_{Y,j}$ of $T_Y$ is normal.

Let $\theta$ be the rational top differential form defining $K_N$.
Define the rational $p_0$-pluricanonical form
\begin{equation}
 \label{eq:slc-cover-generator}
 s=(\pi^*\theta)^{\otimes p_0}\,\pi^*\psi\,w^{-a}.
\end{equation}
Since $\pi^*F=m\Div_Y w$, equations
\eqref{eq:slc-local-relation}--\eqref{eq:slc-cover-crepant} give
\[
 \Div(s)+p_0(T_Y+H_Y)=0.
\]
Thus $s$ generates the degree-$p_0$ divisorial log pluricanonical sheaf
on $Y$.  In particular its sufficiently divisible powers generate the
corresponding invertible log pluricanonical sheaves.

\smallskip
\noindent\emph{Residue comparison on each component.}
Let $T_i$ be the component of $T$ dominated by $T_{Y,j}$, and write
$\rho_j\colon T_{Y,j}\to T_i$ for the induced finite morphism.  At their
generic points the map is \'{e}tale.  The ordinary generic residue
\[
 v_j=\res_{T_{Y,j}}\bigl(s^{p/p_0}\bigr)
\]
is therefore a nonzero rational $p$-pluricanonical form.  The form
$\rho_j^*u_i$ is also nonzero, and the quotient
\begin{equation}
 \label{eq:slc-residue-ratio}
 r_j=\frac{v_j}{\rho_j^*u_i}\in\C(T_{Y,j})^*
\end{equation}
is a rational function.  We prove that $r_j$ is constant by comparing
sufficiently divisible powers.

Choose a multiple $n=bp$ sufficiently divisible for
\eqref{eq:slc-ambient-restriction} and for component adjunction upstairs.
Then $v_j^b$ is a nowhere-vanishing generator of the degree-$n$
adjunction sheaf on $T_{Y,j}$.  The same is true of $\rho_j^*u_i^b$.
For the latter assertion, work near any point of $T$.  By
\eqref{eq:slc-ambient-restriction}, the section $u^b$ is a generator of
the restriction of the invertible sheaf
$\OO_N(n(K_N+T+H))$.  Relative to a local ambient generator it is a unit
on $T$.  Such a unit lifts, after restricting the neighbourhood, to a
unit on $N$.  Hence $u^b$ is the residue restriction of a local ambient
log generator.  Pull that generator back to $Y$.  It remains a log
generator by \eqref{eq:slc-cover-crepant}; its residue on $T_{Y,j}$
generates the upstairs adjunction sheaf.  At the generic point residue
commutes with the \'{e}tale pullback, so this residue is exactly the
rational form $\rho_j^*u_i^b$.  This identifies the rational forms
everywhere they are defined, and proves the assertion at every point
of $T_{Y,j}$.

It follows that $r_j^b$ has zero divisor on $T_{Y,j}$, and hence so does
$r_j$.  A rational function with zero divisor on a normal projective
integral variety is an invertible global function, and therefore a
nonzero complex constant.  We have proved
\begin{equation}
 \label{eq:slc-component-scalar}
 v_j=r_j\rho_j^*u_i,\qquad r_j\in\C^*.
\end{equation}

\smallskip
\noindent\emph{Gluing the scalars and killing the character.}
Suppose two distinct components $T_{Y,j}$ and $T_{Y,j'}$ meet.  Their
intersection contains a stratum of codimension two in $Y$ along which
the pair is generically simple normal crossing.  Its image is a
codimension-two lc centre downstairs.  At the generic point of this
image the map $\pi$ is \'{e}tale and the downstairs pair is simple
normal crossing.  The two downstairs vertical branches are distinct:
\'{e}tale inverse images of one smooth divisor cannot give two
intersecting branches.  Moreover $H$ is absent at this generic
crossing.

At the generic point $\xi$ of the upstairs crossing, choose the two
normal coordinates $x,y$ and a regular ambient $(\dim Y-2)$-form
$\eta$ whose restriction is a nonzero top form on the smooth stratum.
Such choices are available in normal crossing coordinates.  Since
$s^{p/p_0}$ is a log generator, it has the local form
\[
 s^{p/p_0}
   =\varepsilon
      \left(\frac{dx}{x}\wedge\frac{dy}{y}\wedge\eta\right)^{\otimes p},
 \qquad \varepsilon\in\OO_{Y,\xi}^*.
\]
The two next residues differ by $(-1)^p$ and are nonzero because
$\varepsilon$ is a unit.  They therefore agree.  The components of
the tuple $u$ have the same matching, nonzero next residues downstairs:
$u$ is a generator of the slc log pluricanonical line, whose local
description at the node is the even tensor power of the dualizing
line.  Their \'{e}tale pullbacks retain this property.  Taking next
residues in \eqref{eq:slc-component-scalar} now gives $r_j=r_{j'}$.

Since $T_Y$ is connected, its component-incidence graph is connected.
All the constants $r_j$ are consequently equal to one scalar
$r\in\C^*$.  In tuple notation on the normalized components of $T_Y$,
we have
\begin{equation}
 \label{eq:slc-global-scalar}
 \bigl(\res_{T_{Y,j}}(s^{p/p_0})\bigr)_j
     =r\,\pi^*u.
\end{equation}
The tuple $\pi^*u$ is invariant under the covering group, with the
natural action also permuting the components.  On the other hand,
\eqref{eq:slc-cover-generator} shows that for $\zeta\in\mu_m$,
\[
 \zeta^*(s^{p/p_0})=\zeta^{-ap/p_0}s^{p/p_0}.
\]
Residue is natural under this action, including its permutation of the
components.  The nonzero scalar comparison
\eqref{eq:slc-global-scalar} forces this character to be trivial.
Thus $m\mid ap/p_0$.  Since $\gcd(a,m)=1$, we obtain
$m\mid p/p_0$.  Finally
\[
 p\beta=\frac{p}{p_0}\frac{a}{m}\in\Z,
\]
which proves the lemma.
\end{proof}

Apply Lemma~\ref{lem:slc-residue-test} to
\eqref{eq:relative-local-cyclic}, with $\beta=\alpha+t$ and the fixed
coefficient set $\Phi\cup\{1\}$.  It gives a single multiple $p$ of $p_0$
for which \eqref{eq:relative-coefficient-target} holds at every prime
$P$ of the determination $W$.  Consequently $pM_W$ is integral and
Cartier.  Since the moduli b-divisor descends on $W$, the b-divisor
$p\mathbf M$ is b-Cartier.  All constructions used only the dimensions
at most four and the fixed coefficient set.  This completes the proof
of Theorem~\ref{thm:relative-denominators}.

\section{Effective systems recovering the entire base field}
\label{sec:effective}

The exact presentation in the denominator theorem identifies complete
section spaces on the total space and the base. We combine this
identification with effective birationality on a base model.

\begin{proposition}\label{prop:effective-base}
Under the hypotheses of Theorem~\ref{thm:relative-denominators}, suppose in
addition that $D$ is big.  There is a positive integer $m=m(d,\Phi)$ such
that, for every positive multiple $l$ of $m$, the complete divisorial
system
\[
 \left|\left\lfloor l(K_X+B)\right\rfloor\right|
\]
is nonempty and generates precisely the subfield $\C(Z)\subset\C(X)$.
Its rational map is therefore birationally equivalent to $f$ and to the
Iitaka fibration of $K_X+B$.
\end{proposition}

\begin{proof}
Put $L=K_X+B$ and $K=\C(Z)$, and use the exact presentation
\eqref{eq:relative-exact-presentation}. For every positive integer $l$
divisible by $p_0$, we first prove the equality of subspaces of $\C(X)$
\begin{equation}\label{eq:relative-section-spaces}
 H^0\!\left(X,\OO_X(\lfloor lL\rfloor)\right)
   =\psi^{l/p_0}f^*
      H^0\!\left(Z,\OO_Z(\lfloor lD_Z\rfloor)\right).
\end{equation}
Here $f^*$ on the right means pullback of rational functions. For
$v\in K^*$, the exact divisor identity gives
\begin{equation}\label{eq:relative-section-divisors}
 lL+\Div_X\!\left(\psi^{l/p_0}f^*v\right)
    =f^*\!\left(lD_Z+\Div_Z(v)\right).
\end{equation}
If $v$ is a section of the rounded system on $Z$, then
$lD_Z+\Div_Z(v)\geq0$ by the section convention of
Section~\ref{sec:preliminaries}. This divisor is rational Cartier, so
local effective Cartier multiples show that its pullback is effective.
Equation~\eqref{eq:relative-section-divisors} gives the inclusion from
right to left in \eqref{eq:relative-section-spaces}.

Conversely, let $u$ be a nonzero section on the left, and let $X_\eta$
be the generic fibre over $K$. The exact presentation rewrites the
section inequality as
$\Div_X(u/\psi^{l/p_0})+lf^*D_Z\geq0$. Its horizontal valuations give
\[
 \Div_{X_\eta}\!\left(u/\psi^{l/p_0}\right)\geq0.
\]
The generic fibre is normal, so this quotient is regular there. The
contraction condition $f_*\OO_X=\OO_Z$ gives
$H^0(X_\eta,\OO_{X_\eta})=K$. Hence
$u=\psi^{l/p_0}f^*v$ for some $v\in K^*$. By
\eqref{eq:relative-section-divisors}, the source section inequality is
$f^*G\geq0$, where $G=lD_Z+\Div_Z(v)$ is rational Cartier.
Effectivity of $G$ descends to $Z$: near the generic point of any prime
$P\subset Z$, write a Cartier multiple of $G$ as $aP$, with $P$
locally principal. The local Cartier pullback of $P$ has a prime
component dominating $P$, with positive multiplicity $e$, because
$f$ is surjective. Its coefficient in the pullback of the chosen
multiple of $G$ is $ae$. Thus $f^*G\geq0$ forces $a\geq0$.
Doing this for every $P$ proves $G\geq0$, so $v$ is a section on $Z$.
This proves \eqref{eq:relative-section-spaces}. Since $p_0$ clears
$\Phi$, the divisor $lL$ is integral Weil. The argument applies to its
reflexive sheaf even when $lL$ is not Cartier.

We now produce birational sections on the base. Choose a smooth
projective determination $q:W\to Z$ of the nef data which also resolves
the boundary, and write
\[
 K_W+B_W^{\mathrm{cr}}+M_W=q^*D_Z.
\]
Let $A$ be the strict transform of $B_Z$ plus the reduced exceptional
divisor. On nonexceptional primes $A$ agrees with $B_W^{\mathrm{cr}}$;
on exceptional primes its coefficient is one, at least the coefficient
of $B_W^{\mathrm{cr}}$ by generalized log canonicity. Thus
\begin{equation}\label{eq:relative-big-model}
 K_W+A+M_W=q^*D_Z+E_W,
 \qquad E_W=A-B_W^{\mathrm{cr}}\geq0
       \text{ exceptional over }Z.
\end{equation}
Put $\Lambda=\mathcal B\cup\{1\}$. The pair $(W,A)$ is an ordinary
log smooth lc pair with coefficients in the fixed DCC set $\Lambda$,
and $pM_W$ is nef Cartier. Since $D$ is big and $D_Z\sim_\Q D$, the
adjoint on the left of \eqref{eq:relative-big-model} is big.

For each $1\leq k\leq d$, let $b_k=b(\Lambda,k,p)$ be the integer in
Birkar--Zhang's polarized effective birationality theorem
\cite[Theorem~1.3]{BirkarZhang2016}. Its conclusion is that
\[
 \left|\left\lfloor l(K_W+A+M_W)\right\rfloor\right|
\]
is birational for every positive $l$ divisible by $b_k$ when
$\dim W=k$; the floor is taken on the whole adjoint. Set
\[
 m=\operatorname{lcm}(p_0,b_1,\ldots,b_d).
\]
This integer depends only on $d$ and $\Phi$. Fix any positive multiple
$l$ of $m$ and take $k=\dim Z$. For a rational section
$v\in\C(W)=K$ of the displayed birational system, pushforward of its
divisor inequality using \eqref{eq:relative-big-model} gives
\[
 \Div_Z(v)+lD_Z\geq0.
\]
Thus these sections belong to the complete rounded system on $Z$,
and their ratios generate $K$ because the system on $W$ is birational.
The complete system on $Z$ is therefore nonempty and generates $K$.
Equation~\eqref{eq:relative-section-spaces} transfers precisely these
ratios to $X$, since the common factor $\psi^{l/p_0}$ cancels. The
complete source system is nonempty and generates the embedded field
$K\subset\C(X)$ for every such $l$.

The function field of the image of a system's rational map is generated
by its section ratios, so each of these maps is birationally equivalent
to $f$. To compare with the Iitaka fibration, choose for the individual
pair an integer $r>0$ such that $rL$ is Cartier. Degrees divisible by
$\operatorname{lcm}(m,r)$ form a cofinal divisible subsequence of the
Cartier section series of $L$. Such a subsequence has the same field
of section ratios: given $u/v$ in one degree, choose $a\geq1$ so that
the multiplied degree lies in the subsequence and write
$u/v=(uv^{a-1})/v^a$. The field defining the Iitaka
fibration is therefore $K$, and
$\kappa(L)=\operatorname{trdeg}_\C K=\dim Z$. This identifies $f$
with the Iitaka fibration up to birational equivalence. The individual
index $r$ is used only for this comparison and does not enter $m$.
\end{proof}

\section{Integral principal multiples over rationally connected bases}
\label{sec:torsion}

We now bound a common principal multiple of the actual generalized
adjoint, including its denominators.  This will supply uniform
trivializations on rationally connected bases.

\begin{proposition}\label{prop:rc-torsion}
Fix an integer $1\leq d\leq 4$, a DCC set
$I\subset [0,1]\cap\Q$, and a positive integer $p$.
There exists a positive integer $\ell=\ell(d,I,p)$ with the following
property.  Let $Z$ be a normal projective complex variety of dimension $d$
with a rationally connected smooth projective resolution.  Suppose that
$(Z,B+M_Z)$ is generalized klt, that $B\geq 0$ has coefficients in $I$,
and that the b-nef rational b-divisor $\mathbf M$ has $p\mathbf M$
b-Cartier.  If
\[
 D=K_Z+B+M_Z\sim_{\Q}0,
\]
then $\ell D$ is an integral principal divisor.  Consequently, for every
integer $a\geq1$, the divisor $a\ell D$ is principal and
\[
 h^0\bigl(Z,\OO_Z(a\ell D)\bigr)=1.
\]
The same integer can be chosen for all dimensions in any fixed subset of
$\{1,2,3,4\}$.
\end{proposition}

\begin{proof}
There are three steps. Global ACC gives a finite coefficient set and
a uniform positive lower bound for generalized log discrepancies.
Boundedness up to isomorphism in codimension one then bounds torsion
in the divisor class group. Finally we clear the coefficients of the
actual adjoint before applying that torsion bound.

\smallskip
\noindent\emph{Finite coefficients and uniform discrepancies.}
We first construct the small modifications and the one-divisor extractions
that will allow us to use global ACC.  If an exceptional prime divisor $E$
over $Z$ is marked, assume that its generalized log discrepancy is
$a\in(0,1)$.  Choose a sufficiently high projective log resolution
$g:W\to Z$ on which the nef data descend and on which $E$ appears, when
marked.  Write
\[
 K_W+B_W+M_W=g^*D.
\]
The support of $B_W$ is SNC, all its coefficients are less than one, and
$B_W$ need not be effective.

We can also arrange that there is an effective exceptional Cartier divisor
$F$ on $W$ such that $-F$ is $g$-ample and
$\Supp(B_W)\cup\Supp(F)$ is SNC.  To see this, construct a resolution by
blowups with centres over the singular locus of the normal variety $Z$, and
then resolve the nef model, the marked valuation, and the boundary by further
blowups.
The latter can be taken with centres of codimension at least two on smooth
models.  All the resulting exceptional divisors therefore have images of
codimension at least two on $Z$.  At each blowup the negative of its
effective Cartier exceptional divisor is relatively ample.  Adding
sufficiently large multiples of the pullbacks of the divisors chosen at the
preceding stages gives an effective exceptional Cartier divisor whose
negative is ample for the composite morphism.  Additional blowups resolving
the union of the two supports preserve this construction.

Choose a positive rational number $\delta$ sufficiently small that all
coefficients of $B_W+\delta F$ are still less than one and, when $E$ is
marked, $\delta\mult_E F<a$.  For a sufficiently positive ample Cartier
divisor $A$ on $Z$, the divisor
\[
 M_W+g^*A-\delta F
\]
is ample.  Choose a general effective rational divisor
$\Lambda\sim_{\Q}M_W+g^*A-\delta F$ by dividing a general member of a
sufficiently divisible very ample multiple.  Bertini's theorem and the SNC
condition on the union of the supports show that
$(W,B_W+\delta F+\Lambda)$ is sub-klt.
There is a rational principal divisor $P$ on $W$ such that
\[
 \delta F+\Lambda=M_W+g^*A+P.
\]
Since $g$ is birational, $P=g^*(g_*P)$.  Pushing forward gives
\[
 K_Z+B+g_*\Lambda=D+A+g_*P,
\]
and pulling this equality back gives
\[
 g^*(K_Z+B+g_*\Lambda)
       =K_W+B_W+\delta F+\Lambda.
\]
Thus, setting $\Xi=B+g_*\Lambda$, the pair $(Z,\Xi)$ is an ordinary
effective rational klt pair.  If $E$ is marked, the general divisor
$\Lambda$ does not contain $E$, so its ordinary log discrepancy is
\[
 a(E,Z,\Xi)=a-\delta\mult_E F.
\]
This lies in $(0,1)$ by the choice of $\delta$ and the original bound $a<1$.

Take a log resolution of $(Z,\Xi)$ dominating $W$, so that it exhibits
$E$ when marked.  Apply the extraction corollary
\cite[Corollary~1.4.3 and its proof]{BCHM2010} with the prescribed set
$\{E\}$, or with the empty set when no valuation is marked.  The klt pair
$(Z,\Xi)$ satisfies the lc hypothesis, and $\Xi$ itself can be used as the
auxiliary klt boundary.  When a divisor is marked, its discrepancy is less
than one, so the extra condition for selected discrepancy-one valuations is
vacuous.  We obtain
a birational morphism
\[
 h:Z^+\longrightarrow Z
\]
with $Z^+$ $\Q$-factorial and with exactly the prescribed exceptional prime
divisors.  The construction in the cited proof is a log terminal model over
$Z$, hence $h$ is projective.  In the empty case it is a small
$\Q$-factorial modification.

Restore the original generalized data on $Z^+$ by setting
\[
 K_{Z^+}+B^++M_{Z^+}=h^*D,
\]
with the same nef b-divisor $\mathbf M$.  The resulting pair is generalized
klt and its adjoint is rationally linearly trivial.  Moreover, $B^+$ is
effective: it is the strict transform of $B$, together with $(1-a)E$ in
the extraction case.  If a marked divisor is already a prime on $Z$, the
small modification given by the empty-list construction suffices.

We next apply global ACC to these effective generalized pairs, whenever
their boundary coefficients belong to a fixed DCC set.  There are two
cases.  If the nef trace is not numerically trivial on a model determining
$\mathbf M$, then on a common higher model the nef part has the form
\[
 M_W=\frac1p(pM_W),
\]
where $pM_W$ is a nef Cartier divisor that is not numerically trivial.
The generalized global ACC theorem
\cite[Theorem~1.6]{BirkarZhang2016} applies: the dimension is fixed, the
pair is projective generalized lc, its total adjoint is numerically
trivial, the boundary coefficients belong to the chosen DCC set, and the
single nef weight is $1/p$.  We adjoin this fixed weight to that DCC set.
It follows that the possible boundary coefficients belong to a finite
set.

If instead $M_W\equiv0$, take a common higher model
$\pi:\widetilde W\to Z^+$ on which the nef data descend.  Since $Z^+$ is
$\Q$-factorial, the trace $M_{Z^+}$ is $\Q$-Cartier.  The divisor
\[
 M_{\widetilde W}-\pi^*M_{Z^+}
\]
is exceptional and numerically trivial over $Z^+$.  Applying the
negativity lemma to it and to its negative shows that it is zero.
Consequently $M_{Z^+}\equiv0$: every curve on $Z^+$ is dominated by a curve
on $\widetilde W$, and we use the projection formula.  Since the nef data
are pulled back from $Z^+$, the generalized discrepancies are the ordinary
discrepancies of $(Z^+,B^+)$.  This is an ordinary klt pair with
$K_{Z^+}+B^+\equiv0$, so ordinary global ACC
\cite[Theorem~1.5]{HMX2014} again gives a finite set of boundary
coefficients.

Applying this discussion to the small modification with no marked divisor
first shows that the coefficients of $B$ lie in a fixed finite rational
set $I_0\subset I$.  It also shows that the original generalized pairs
are uniformly generalized $\varepsilon$-lc for some
$\varepsilon=\varepsilon(d,I,p)>0$.  Indeed, otherwise we could choose a
sequence of such pairs and marked valuations whose generalized log
discrepancies satisfy
\[
 1>a_1>a_2>\cdots>0,\qquad a_i\longrightarrow0.
\]
Perform the extraction just constructed when the marked valuation is
exceptional, and otherwise take the small modification.  The boundary
coefficients of the resulting pairs belong to
\[
 I\ \cup\ \{1-a_i:i\geq1\},
\]
which is a DCC set because the added sequence is strictly increasing.
Global ACC would force the coefficients $1-a_i$ to belong to a finite
set, a contradiction.

\smallskip
\noindent\emph{Bounded topology and class-group torsion.}
We have obtained the uniform singularity bound without choosing any
bound for the class group. We now turn this geometric information into
a finite topological list.

Birkar's boundedness theorem
\cite[Theorem~1.7]{BirkarBoundedCY} now applies to the original varieties
$Z$.  Its hypotheses are
projectivity, fixed dimension, generalized $\varepsilon$-lc singularities,
effective boundary, a real-linearly trivial generalized adjoint, and
rational connectedness.  The singularity bound was just proved, the
adjoint is even rationally linearly trivial, and rational connectedness follows
from that of the given smooth projective resolution.  Thus, for every $Z$,
the theorem gives a projective variety $V_Z$, chosen from one bounded set,
that is isomorphic to $Z$ in codimension one.  Replacing each $V_Z$ by its
finite normalization preserves this property: normalization is an
isomorphism over the common normal open subset, and the complement still
has codimension at least two.  These normalizations still form a bounded set,
as in \cite[Section~9.5]{BirkarBoundedCY}.  Hence there is a projective
morphism $\mathcal V\to S$, with $S$ of finite type over $\C$, such that for
every $Z$ some fibre $V_s:=\mathcal V_s$, $s\in S(\C)$, is normal and
isomorphic to $Z$ in codimension one.

We explain explicitly how this boundedness controls torsion.  Put
$U=Z_{\mathrm{reg}}$.  The groups
\[
 H_1\bigl(U(\C),\Z\bigr)
\]
are finitely generated and take only finitely many isomorphism types.
First, this group is invariant under isomorphism in codimension one of
normal varieties.  The common open subset, intersected with the regular
loci, is obtained from each regular locus by removing a closed subset of
complex codimension at least two.  Removing such a subset from a smooth
complex variety does not change its fundamental group: choose a smooth
stratification of the removed subset and perturb paths and homotopies,
relative to their prescribed boundaries, to be transverse to its strata.
All these strata have real codimension at least four, so the perturbed
paths and two-dimensional homotopies avoid them.  The assertion for
$H_1$ follows by abelianization.

Second, we prove the assertion for the selected fibres $V_s$.  Replace $S$
by its reduction and partition it into finitely many locally closed strata
over which the family is flat.  Cover these strata by finitely many affine
opens.  On any one of these opens, still denoted $S$, let
\[
 \mathcal U=\operatorname{Sm}(\mathcal V/S)
\]
be the full relative smooth locus.  Flatness and finite presentation give
\[
 \mathcal U_s=\operatorname{Sm}(\mathcal V_s/\C)=(V_s)_{\mathrm{reg}}
\]
for each selected normal fibre, since regularity and smoothness agree over
$\C$.  Embed $\mathcal V$ in $\PP^N_{\C}\times S$.  Identify $\PP^N(\C)$
with the real algebraic set of Hermitian matrices $H$ satisfying $H^2=H$
and $\tr H=1$, by $[z]\mapsto zz^*/(z^*z)$, where $z^*$ is conjugate
transpose.  With real and imaginary coordinates on the affine $S$, the
entire map
\[
 \mathcal U(\C)\longrightarrow S(\C)
\]
is then a continuous semialgebraic map between affine semialgebraic spaces.
The semialgebraic triviality theorem of Delfs and Knebusch
\cite[Theorem~6.4]{DelfsKnebusch1982} gives a finite partition of $S(\C)$
into semialgebraic pieces over each of which the whole inverse image is a
product.  Thus the selected regular loci have only finitely many
homeomorphism types.  Each semialgebraic fibre has the homotopy type of a
finite simplicial complex by
\cite[Theorem~2.1 and Propositions~2.4--2.5]{DelfsKnebusch1982}.
The asserted finite list of finitely generated groups $H_1$ follows.

For each individual $Z$, the group $\Cl(Z)$ is also finitely generated.
Indeed, let $r:Y\to Z$ be a smooth projective rationally connected
resolution.  Weil divisor pushforward gives a surjection
\[
 \Pic(Y)\longrightarrow\Cl(Z):
\]
strict transforms give surjectivity on divisors, and principal divisors
push forward to principal divisors.  The Albanese map of $Y$ is constant
on every rational curve and therefore constant on $Y$, since rational
curves connect general points.  Consequently $\Pic^0(Y)=0$, and
N\'eron--Severi finiteness makes $\Pic(Y)$ finitely generated.

Normality and projectivity give
$\Gamma(U,\OO_U^*)=\C^*$, and restriction gives
$\Pic(U)=\Cl(Z)$ because $U$ is regular and contains every
codimension-one point.  The $n$th-power map on $\C^*$ is surjective.
For every $n\geq1$, the \'etale Kummer sequence
\cite[Tag~03PK]{StacksProject} gives the first isomorphism below.
Riemann existence for finite \'etale covers
\cite[Expos\'e~XII, Th\'eor\`eme~5.1 and Corollaire~5.2]{SGA1} gives the
second by classification of finite abelian torsors and abelianization:
\begin{equation}\label{eq:rc-kummer}
 \Cl(Z)[n]
 \simeq H^1_{\mathrm{\acute et}}(U,\mu_n)
 \simeq
 \Hom\bigl(H_1(U(\C),\Z),\mu_n(\C)\bigr).
\end{equation}
For a fixed $Z$, the orders of the groups on the left are bounded as $n$
varies, because $\Cl(Z)$ is finitely generated.  If $H_1(U(\C),\Z)$ had
a free summand of positive rank, the orders of the groups on the right
would grow without bound.  Hence $H_1(U(\C),\Z)$ is finite.  Only finitely
many such finite groups occur by the preceding paragraph, so there is a
uniform integer $T$ divisible by all their exponents.  Every torsion class
$c\in\Cl(Z)$ is killed by $T$: take $n$ equal to its order and apply
\eqref{eq:rc-kummer}, whose right-hand side has exponent dividing $T$.

\smallskip
\noindent\emph{An integral principal multiple.}
Finally, $pM_Z$ is an integral Weil divisor, being the pushforward of a
Cartier divisor on a model determining the nef data.  Choose a uniform
integer $q$ divisible by $p$ and by all denominators in the finite set
$I_0$.  Then $qD$ is an integral Weil divisor.  The relation
$D\sim_{\Q}0$ makes the class $[qD]\in\Cl(Z)$ torsion: if
$nD=\Div(v)$ for some $n\geq1$ and $v\in\C(Z)^*$, then
$n(qD)=\Div(v^q)$.  Hence
$TqD$ is principal.  Taking $\ell=Tq$ proves the proposition for dimension
$d$.  Taking a common multiple of the finitely many resulting integers
handles the stated range of dimensions.  All positive multiples of
$\ell D$ are principal, and projectivity and integrality of $Z$ give the
last assertion about sections.
\end{proof}

\begin{corollary}\label{cor:relative-rc-torsion}
Fix a finite rational coefficient set $\Phi\subset[0,1]\cap\Q$.
There is a positive integer $r=r(\Phi)$ such that the following holds.
Let $(X,B)$ be a projective klt fourfold pair with effective boundary
having coefficients in $\Phi$, and suppose that
$K_X+B\sim_{\Q}0$.  If there is a contraction $f:X\to Z$ with
$\dim Z>0$ and with a rationally connected smooth projective resolution
of $Z$, then $r(K_X+B)$ is an integral principal divisor. This includes
$B=0$, and every positive multiple is also principal.
\end{corollary}

\begin{proof}
Apply Theorem~\ref{thm:relative-denominators} to $f$ with the zero divisor
as the rational pullback class.  It gives uniformly bounded integers
$p_0\mid p$, a rational function $\psi\in\C(X)^*$, and an exact divisor
identity
\begin{equation}\label{eq:rc-exact-relative}
 K_X+B+\frac1{p_0}\Div(\psi)=f^*D_Z,
 \qquad D_Z=K_Z+B_Z+M_Z,
\end{equation}
where $(Z,B_Z+M_Z)$ is generalized klt, $B_Z$ is effective with
coefficients in a fixed rational DCC set, and $p\mathbf M$ is b-Cartier
with nef data.  The divisor $D_Z$ is rationally linearly trivial.  This
also follows directly from $f^*D_Z\sim_{\Q}0$: choose a positive
multiple making $D_Z$ Cartier and its pullback linearly trivial, and use
$f_*\OO_X=\OO_Z$ and the projection formula to trivialize that multiple
on $Z$.

Proposition~\ref{prop:rc-torsion} gives a uniform integer $\ell$ and
$v\in\C(Z)^*$ such that $\ell D_Z=\Div(v)$.  Choose $r$ divisible by both
$\ell$ and $p$, so that $p_0\mid r$.  Multiplying
\eqref{eq:rc-exact-relative} by $r$ gives
\[
 r(K_X+B)
 =\Div\!\left((v\circ f)^{r/\ell}\psi^{-r/p_0}\right).
\]
Both exponents are integers, so this is an equality to the divisor of an
actual rational function and includes the integrality of the adjoint
multiple.  All constants depend only on $\Phi$, since the possible
dimensions of $Z$ form a finite
set and the denominator theorem supplies uniform $p_0,p$ and a fixed
DCC set.
\end{proof}

\begin{thebibliography}{99}

\bibitem{AmbroBoundary2004}
F.~Ambro, \emph{Shokurov's boundary property},
J. Differential Geom. \textbf{67} (2004), no.~2, 229--255.
\href{https://doi.org/10.4310/jdg/1102536201}{doi:10.4310/jdg/1102536201}.

\bibitem{AmbroModuli2005}
F.~Ambro, \emph{The moduli b-divisor of an lc-trivial fibration},
Compos. Math. \textbf{141} (2005), no.~2, 385--403.
\href{https://doi.org/10.1112/S0010437X04001071}{doi:10.1112/S0010437X04001071}.

\bibitem{BakkerFilipazziMauriTsimerman2025}
B.~Bakker, S.~Filipazzi, M.~Mauri, and J.~Tsimerman,
\emph{Baily--Borel compactifications of period images and the b-semiampleness conjecture},
preprint, \href{https://arxiv.org/abs/2508.19215v2}{arXiv:2508.19215v2},
18 December 2025.

\bibitem{BCHM2010}
C.~Birkar, P.~Cascini, C.~D.~Hacon, and J.~McKernan,
\emph{Existence of minimal models for varieties of log general type},
J. Amer. Math. Soc. \textbf{23} (2010), no.~2, 405--468.
\href{https://doi.org/10.1090/S0894-0347-09-00649-3}{doi:10.1090/S0894-0347-09-00649-3}.
\href{https://arxiv.org/abs/math/0610203v2}{arXiv:math/0610203v2}.

\bibitem{BirkarBoundedCY}
C.~Birkar,
\emph{Singularities on Fano fibrations and beyond},
preprint, \href{https://arxiv.org/abs/2305.18770v2}{arXiv:2305.18770v2}
(2025).

\bibitem{BirkarZhang2016}
C.~Birkar and D.-Q.~Zhang,
\emph{Effectivity of Iitaka fibrations and pluricanonical systems of polarized pairs},
Publ. Math. Inst. Hautes \'Etudes Sci. \textbf{123} (2016), 283--331.
\href{https://doi.org/10.1007/s10240-016-0080-x}{doi:10.1007/s10240-016-0080-x}.

\bibitem{ChenTsakanikas2023}
G.~Chen and N.~Tsakanikas,
\emph{On the termination of flips for log canonical generalized pairs},
Acta Math. Sin. (Engl. Ser.) \textbf{39} (2023), 967--994.
\href{https://doi.org/10.1007/s10114-023-0116-3}{doi:10.1007/s10114-023-0116-3}.
Preprint: \href{https://arxiv.org/abs/2011.02236v2}{arXiv:2011.02236v2}.

\bibitem{DelfsKnebusch1982}
H.~Delfs and M.~Knebusch,
\emph{On the homology of algebraic varieties over real closed fields},
J. reine angew. Math. \textbf{335} (1982), 122--163.
\href{https://doi.org/10.1515/crll.1982.335.122}{doi:10.1515/crll.1982.335.122}.
\href{https://epub.uni-regensburg.de/12717/1/ubr05118_ocr.pdf}{Repository copy}.

\bibitem{Floris2013}
E.~Floris, \emph{Bounds on the denominators in the canonical bundle formula},
Ann. Inst. Fourier (Grenoble) \textbf{63} (2013), no.~5, 1951--1969.
\href{https://doi.org/10.5802/aif.2819}{doi:10.5802/aif.2819}.

\bibitem{Fujino2000}
O.~Fujino,
\emph{Abundance theorem for semi log canonical threefolds},
Duke Math. J. \textbf{102} (2000), no.~3, 513--532.
\href{https://doi.org/10.1215/S0012-7094-00-10237-2}{doi:10.1215/S0012-7094-00-10237-2}.
\href{https://www.math.kyoto-u.ac.jp/~fujino/Abundance.pdf}{Author's copy of the published article}.

\bibitem{FujinoGongyoModuli}
O.~Fujino and Y.~Gongyo,
\emph{On the moduli b-divisors of lc-trivial fibrations},
Ann. Inst. Fourier (Grenoble) \textbf{64} (2014), no.~4, 1721--1735.
\href{https://doi.org/10.5802/aif.2894}{doi:10.5802/aif.2894}.

\bibitem{FujinoMori2000}
O.~Fujino and S.~Mori, \emph{A canonical bundle formula},
J. Differential Geom. \textbf{56} (2000), no.~1, 167--188.
\href{https://doi.org/10.4310/jdg/1090347529}{doi:10.4310/jdg/1090347529}.

\bibitem{Gongyo2013}
Y.~Gongyo,
\emph{Abundance theorem for numerically trivial log canonical divisors of semi-log canonical pairs},
J. Algebraic Geom. \textbf{22} (2013), no.~3, 549--564.
\href{https://doi.org/10.1090/S1056-3911-2012-00593-1}{doi:10.1090/S1056-3911-2012-00593-1}.
\href{https://www.ms.u-tokyo.ac.jp/~gongyo/papers/abundance_sub7.pdf}{Author version}.

\bibitem{SGA1}
A.~Grothendieck (dir.),
\emph{Rev\^etements \'etales et groupe fondamental (SGA~1)},
S\'eminaire de G\'eom\'etrie Alg\'ebrique du Bois Marie 1960--61,
augment\'e de deux expos\'es de M.~Raynaud,
Documents Math\'ematiques \textbf{3},
Soci\'et\'e Math\'ematique de France, Paris, 2003.
\href{https://arxiv.org/abs/math/0206203}{arXiv:math/0206203}.

\bibitem{HMX2014}
C.~D.~Hacon, J.~McKernan, and C.~Xu,
\emph{ACC for log canonical thresholds},
Ann. of Math. (2) \textbf{180} (2014), no.~2, 523--571.
\href{https://doi.org/10.4007/annals.2014.180.2.3}{doi:10.4007/annals.2014.180.2.3}.

\bibitem{JiangLiu2021}
C.~Jiang and H.~Liu,
\emph{Boundedness of log pluricanonical representations of log Calabi--Yau
pairs in dimension 2},
Algebra Number Theory \textbf{15} (2021), no.~2, 545--567.
\href{https://doi.org/10.2140/ant.2021.15.547}{doi:10.2140/ant.2021.15.547}.
Preprint: \href{https://arxiv.org/abs/2002.11928v1}{arXiv:2002.11928v1}.

\bibitem{Kollar2013}
J.~Koll\'ar, with the collaboration of S.~Kov\'acs,
\emph{Singularities of the minimal model program},
Cambridge Tracts in Mathematics, vol.~200,
Cambridge University Press, Cambridge, 2013.
\href{https://doi.org/10.1017/CBO9781139547895}{doi:10.1017/CBO9781139547895}.

\bibitem{KollarFamilies2023}
J.~Koll\'ar, with the collaboration of K.~Altmann and S.~J.~Kov\'acs,
\emph{Families of varieties of general type},
Cambridge Tracts in Mathematics, vol.~231,
Cambridge University Press, Cambridge, 2023.
\href{https://doi.org/10.1017/9781009346115}{doi:10.1017/9781009346115}.
\href{https://web.math.princeton.edu/~kollar/FromMyHomePage/modbook-final.pdf}{Author version dated 8 August 2022}.

\bibitem{KollarMori1998}
J.~Koll\'ar and S.~Mori,
\emph{Birational Geometry of Algebraic Varieties},
Cambridge Tracts in Mathematics \textbf{134},
Cambridge University Press, Cambridge, 1998.
\href{https://doi.org/10.1017/CBO9780511662560}{doi:10.1017/CBO9780511662560}.

\bibitem{StacksProject}
The Stacks Project Authors,
\emph{The Stacks Project},
Lemma~58.21.4, Tag
\href{https://stacks.math.columbia.edu/tag/0BMB}{0BMB}
(purity of the branch locus), and Tag
\href{https://stacks.math.columbia.edu/tag/03PK}{03PK}
(Kummer theory).

\bibitem{TodorovXu2009}
G.~Todorov and C.~Xu,
\emph{Effectiveness of the log Iitaka fibration for 3-folds and 4-folds},
Algebra Number Theory \textbf{3} (2009), no.~6, 697--710.
\href{https://doi.org/10.2140/ant.2009.3.697}{doi:10.2140/ant.2009.3.697}.

\end{thebibliography}

\end{document}
